\section{Optical sensitivity and source-weighted exit flux}
\label{sec:v79-optical-flux}
The two terms in \eqref{eq:v79-separated-endpoint} have distinct
geometric mechanisms. We first give a positive, collision-uniform
sensitivity formula for the optical term, and then an exact boundary
formula for the initial source mass which exits. Neither formula
identifies a surface insertion with a bounded volume multiplier.

\subsection{A positive optical polynomial}
For a regular word put $t_j=\ell_j^{-1}$, $V_i=2/(Rc_i)$ and
\begin{equation}\label{eq:v79-optical-matrix}
 A_{ii}=t_{i-1}+t_i+V_i,\qquad A_{i,i+1}=A_{i+1,i}=-t_i.
\end{equation}
The matrix $A$ is the incidence-normalized optical matrix; the actual
internal arclength Hessian is $D_cAD_c$. In the Neumann estimate it
is $A=D_A(I-P)$, where $D_A$ is its diagonal and the row sums of $P$
are at most $q_0$. These two diagonal normalizations are different.

Define $y_0=0$, $f_0=1$, $y_1=\ell_0$, and
\begin{equation}\label{eq:v79-positive-recursion}
 f_i=f_{i-1}+V_i y_i,\qquad y_{i+1}=y_i+\ell_i f_i\quad(1\le i<m).
\end{equation}
The continuant and corner Schur-complement identities give
\begin{equation}\label{eq:v79-optical-denominator}
 y_m=\left(\prod_{j=0}^{m-1}\ell_j\right)\det A,
 \qquad \omega=\frac{R c_0c_m}{y_m}.
\end{equation}
For $m=1$ the determinant is empty and $y_1=\ell_0$.
Indeed eliminating $f_i$ gives
$y_{i+1}=(1+\ell_i/\ell_{i-1}+\ell_iV_i)y_i
-(\ell_i/\ell_{i-1})y_{i-1}$, the continuant recurrence with the same
initial values. In the mixed endpoint Schur complement all internal
incidence squares cancel against $\det(D_cAD_c)$, leaving
$|F_{uv}|=c_0c_m\prod t_j/\det A$ in endpoint arclengths. The change
to endpoint angles gives $\omega=R|F_{uv}|$ as stated.

\begin{proposition}[Optical logarithmic sensitivities]
\label{prop:v79-optical-sensitivity}
The polynomial $y_m$ has nonnegative coefficients and is of degree
at most one in each of the separate variables $\ell_j,V_i$. Thus
\begin{equation}\label{eq:v79-sensitivity-weights}
 a_i=V_i\partial_{V_i}\log y_m\in[0,1],\qquad
 b_j=\ell_j\partial_{\ell_j}\log y_m\in[0,1].
\end{equation}
At fixed $R$ its differential on any actual regular word is
\begin{equation}\label{eq:v79-log-optical}
 d\log\omega=d\log c_0+d\log c_m
       +\sum_{i=1}^{m-1}a_i\,d\log c_i
       -\sum_{j=0}^{m-1}b_j\,d\log\ell_j.
\end{equation}
In endpoint-angle coordinates, for a fixed uniform $C_O$,
\begin{equation}\label{eq:v79-optical-gradient}
 \|\nabla_q\log\omega\|
       \le C_O\left(1+\sum_{i=1}^{m-1}\frac{q_0^{d_i}}{c_i}\right).
\end{equation}
No third derivative or sum of absolute word multiplicities occurs
in this bound.
\end{proposition}
\begin{proof}
Each new $V_i$ enters the positive recursion once, and each new
$\ell_i$ enters once after that step. Later operations are linear in
all previous variables. Induction proves positivity and separate
affineness. For a positive separately affine polynomial $Y=A+Bz$,
$z\partial_z\log Y=Bz/(A+Bz)\in[0,1]$. Applying this identity to each
variable gives \eqref{eq:v79-sensitivity-weights}; differentiating
\eqref{eq:v79-optical-denominator} gives \eqref{eq:v79-log-optical}.

The endpoint incidences are bounded below by $\sqrt{19}/10$.
Theorem~\ref{thm:v78-margins} bounds $\|D_qc_i\|$ by $C_Dq_0^{d_i}$.
Also $|d\ell_j|\le|dQ_j|+|dQ_{j+1}|$ and $\ell_j\ge\tau_-$, so
$\|D_q\log\ell_j\|\le2C_Qq_0^{d'_j}/\tau_-$. The sum of these
flight bounds is uniform in $m$. Inserting the sensitivity weights
proves \eqref{eq:v79-optical-gradient}. The inverse-incidence terms
are retained; positivity does not make their complete source sum finite.
\end{proof}

\subsection{A quantified regular part of the optical loss}
For $0<\delta\le1$ let $S_h^\delta(t)\subset S_h(t)$ consist of
survivors satisfying, for every internal contact and every $0\le s\le h$,
\begin{equation}\label{eq:v79-incidence-corridor}
               c_i(\Phi_sx)\ge\delta q_0^{d_i/2}.
\end{equation}
Write $O_h^{\varepsilon,\mathrm{reg}(\delta)}$ for
\eqref{eq:v79-optical-loss} restricted to this set, and
$O_h^{\varepsilon,\mathrm{thin}(\delta)}$ for its nonnegative complement.
This is a decomposition of the original weighted source, not a new
guard or a renormalized probability law.

\begin{theorem}[Uniform regular optical distortion]
\label{thm:v79-regular-optical}
There are uniform positive $c_1,C_1,C_2$ such that, whenever
$0<r\le c_1\delta^2$,
\begin{align}
 \limsup_m\sup_{R,n,k}m^2
 \|\bar O_r^{\varepsilon,\mathrm{reg}(\delta)}\|_\infty
                         &\le C_1\delta^{-1}\sqrt r,
                                      \label{eq:v79-regular-optical}\\
 \bar O_r^{\varepsilon,\mathrm{thin}(\delta)}(t)
                         &\le b^{C_2\delta}(t)
                       \quad\text{for almost every }t.
                                      \label{eq:v79-thin-optical}
\end{align}
The first estimate uses only the stated fixed-width analytic input.
The second is a pointwise source domination and is not a bound for
the normalized essential height of its right side.
\end{theorem}
\begin{proof}
On the corridor, the sum in \eqref{eq:v79-optical-gradient} is at most
$2\delta^{-1}/(1-\sqrt{q_0})$. The forward endpoint-curve length is
at most $C_X\sqrt h$. Hence $|\log\eta_h|\le C\delta^{-1}\sqrt h$.
Choose $c_1$ so that this upper bound is at most one for $h\le r$.
Then $[1-\eta_h]_+\le C\delta^{-1}\sqrt h\,\eta_h$.
As in the proof of Theorem~\ref{thm:v79-volume-factorization},
$\eta_h\le(\Lambda_\circ/\lambda_\circ)J_h$ and target coarea gives
\[
 O_h^{\varepsilon,\mathrm{reg}(\delta)}(t)
                     \le C\delta^{-1}\sqrt h\,p(t+h).
\]
Averaging first proves \eqref{eq:v79-regular-optical}.

If the corridor fails at time $s\le h\le r$, the margin motion bound
implies at the starting point
\[
 c_i(x)\le\delta q_0^{d_i/2}+C_DC_Xq_0^{d_i}\sqrt r
                          \le C_2\delta q_0^{d_i/2}.
\]
Since $q_0^{d_i/2}\le q_0^{d_i/4}$, the corresponding factor of the
original guard at auxiliary parameter $C_2\delta$ is zero. Thus
$W^{C_2\delta}(x)=1$ on the entire thin starting set. Its optical
loss is bounded by its original weighted source, and $W^\varepsilon\le1$.
This proves \eqref{eq:v79-thin-optical} before averaging, hence afterwards.
There are no internal contacts when $m=1$ and the thin set is empty.
\end{proof}

The theorem quantitatively controls the regular optical contribution
on the full exact event. It does not allow replacement of
$b^{C_2\delta}$ by its fixed-width average at a starting roof. That
replacement would be exactly the missing pointwise step.

\subsection{Exit mass in the correct flux convention}
The initial mass of trajectories which exit need not equal the
physical flux measured at their exit. Compression between the starting
roof and the exit must be retained. We state the exact formula on
regular subdomains and recover the full stopped event by exhaustion.

Let $D$ be a relatively compact regular domain in the exact event,
away from $g=0$, with piecewise smooth boundary. We take the boundary
transverse to $X$ except for pieces and lower-dimensional strata whose
images in the angular coordinate have zero length. Such domains are
obtained from finite flow boxes; increasing exhaustions can be formed
by refining the boxes. Let $\Gamma_D^+$ denote its outward boundary,
$N$ its outward unit normal in the collision coordinates, and
$\tau_D^-(y)$ the first backward lifetime from $y$ inside $D$.
Use the physical flux measure
\begin{equation}\label{eq:v79-boundary-flux}
 d\zeta_D(y)=\frac{[X(y)\cdot N(y)]_+}{4\pi c_*}\,d\mathcal H^1(y).
\end{equation}

\begin{proposition}[First-exit flux with the initial source weight]
\label{prop:v79-source-exit-flux}
Let $E_{D,h}^\varepsilon(t)$ be the weighted initial mass on $L=t$
starting in $D$ and exiting $D$ within $h$. For almost every $(t,h)$,
\begin{equation}\label{eq:v79-exit-flux}
 E_{D,h}^\varepsilon(t)=
 \int_{\Gamma_D^+}
 \one_{\{0<L(y)-t\le h,\ L(y)-t<\tau_D^-(y)\}}
 \frac{W^\varepsilon(x_{t,y})}
      {J_{L(y)-t}(x_{t,y})}\,d\zeta_D(y),
 \quad x_{t,y}=\Phi_{t-L(y)}y.
\end{equation}
For an increasing regular exhaustion $D_\ell$ of the full regular
exact event, its original true-exit mass satisfies
\begin{equation}\label{eq:v79-exhausted-exit}
 E_h^\varepsilon(t)=\lim_{\ell\to\infty}
 \left\{\int_{\{L=t\}\setminus D_\ell}W^\varepsilon\,d\mu_t
                                  +E_{D_\ell,h}^\varepsilon(t)\right\}.
\end{equation}
The bracketed nonnegative quantities decrease to the left side.
The limit is at fixed collision count; it is not interchanged with
a normalized collision limsup.
\end{proposition}
\begin{proof}
In a flow box use $(u,\theta)=(L,\vartheta)$, so that $X=\partial_u$.
The physical source has area density
$a(u,\theta)=\omega/(4\pi c_*\chi)$ relative to $du\,d\theta$.
On a transverse outward boundary $u=b(\theta)$, its flux is
$a(b(\theta),\theta)|d\theta|$. At the starting roof its coarea
mass is $a(t,\theta)|d\theta|$, and
\[
 J_{b(\theta)-t}(x_{t,y})
                        =a(b(\theta),\theta)/a(t,\theta).
\]
Multiplication by the initial weight proves \eqref{eq:v79-exit-flux}
on the box. The backward lifetime selects the first forward exit
rather than a later re-entry. Uniqueness of the flow makes the source
pieces disjoint; a measurable refinement of the boxes sums them
without multiplicity. The excluded boundary strata have zero angular
measure and hence zero source-fiber measure. Coarea supplies the joint
almost-everywhere statement.

For each starting point, declare the $D_\ell$-lifetime zero when it
starts outside $D_\ell$. These lifetimes increase to the lifetime in
the full event: every compact surviving flow segment eventually lies
in some $D_\ell$. The indicators of lifetime at most $h$ therefore
decrease to the corresponding full-event indicator. Dominated
convergence on each finite coarea fiber proves
\eqref{eq:v79-exhausted-exit}. The exceptional $(t,h)$ set can be
chosen jointly by disintegration. Chart boundaries in the exhaustion
are not declared physical; their effects disappear only through this
source-level limiting identity.
\end{proof}

The inverse Jacobian in \eqref{eq:v79-exit-flux} is essential. Along
a limiting grazing boundary the physical density may tend to zero
while this inverse factor diverges. Vanishing unweighted exit flux
therefore does not make the original initial exit mass zero. The
formula retains grazing, competing-hit and section exits. If several
margins first vanish simultaneously, a fixed ordering partitions the
initial source into disjoint exit types; their sum is $E_h^\varepsilon$.
The internal occupation-current identity in the supplement is signed.
It has not been used to subtract any of these positive exit masses.
